\section{Proposed Method: Results}
\label{sec:results-proposed}

This section focuses on our \textbf{asymmetric hybrids}
(Section~\ref{sec:method}) relative to the dual-InfoNCE starting point.
The full published-loss ablation on the same 3M encoder at 200 epochs is in
Section~\ref{sec:ablation}.

\subsection{Dual InfoNCE baseline}
\label{sec:results-dual}

Replacing cosine JEPA with batch InfoNCE on both the mask and augmentation
paths (\texttt{latent\_infonce\_jepa\_augment}) yields strong instance
discrimination: k-NN@20 reaches $\sim$0.21--0.23 by epoch~90, but linear
probe accuracy remains near $\sim$0.10.
Embedding visualizations (PCA, k-NN neighborhood panels in our analysis
scripts) show class-mixed local neighborhoods---consistent with high k-NN
driven by instance structure rather than semantic clusters.

\subsection{Asymmetric hybrid results (@100 epochs)}
\label{sec:results-hybrid}

Table~\ref{tab:hybrid} summarizes the two proposed second-term designs at
100 epochs (500-epoch extension in progress for InfoNCE+SIGReg).

\begin{table}[t]
  \centering
  \caption{Proposed asymmetric hybrids vs.\ dual InfoNCE (SmallResNet,
  CIFAR-100).}
  \label{tab:hybrid}
  \begin{tabular}{lccc}
    \toprule
    Method & k-NN@20 & Linear & Epoch \\
    \midrule
    Dual InfoNCE (baseline)     & 0.21--0.23 & 0.103 & 90 \\
    InfoNCE + MSE + var + cov     & 0.190 & 0.147 & 30 \\
    InfoNCE + MSE + SIGReg        & 0.179 & \textbf{0.171} & 100 \\
    InfoNCE + MSE + SIGReg (cont.)& --- & --- & 500 \\
    \bottomrule
  \end{tabular}
\end{table}

\paragraph{Effect of the second term.}
Swapping augmentation InfoNCE for spread regularizers increases linear
accuracy by $43$--$66\%$ relative to dual InfoNCE while reducing k-NN by
roughly $0.02$--$0.04$.
Decomposed variance--covariance (unit weights) preserves more k-NN; SIGReg
(unit weight) pushes linear accuracy further at a larger k-NN cost.

\subsection{Loss diagnostics}
\label{sec:results-diagnostics}

For InfoNCE+SIGReg at epoch~100, the total loss is dominated by
$\mathcal{L}_{\mathrm{sig}}$ ($\sim$2.33 of $\sim$2.84), with
$\mathcal{L}_{\mathrm{mask}}\approx 0.24$ and $\mathcal{L}_{\mathrm{inv}}
\approx 0.27$.
Mask-path InfoNCE remains well optimized; aug--clean MSE is looser than in
the MVC hybrid ($\mathcal{L}_{\mathrm{inv}}\approx 0.03$ at epoch~30),
indicating SIGReg-driven spread without tight augmentation invariance.

\begin{table}[t]
  \centering
  \caption{InfoNCE + MSE + SIGReg loss components at epoch~100.}
  \label{tab:loss}
  \begin{tabular}{lrr}
    \toprule
    Term & Value & Weight \\
    \midrule
    InfoNCE (mask) & 0.236 & 1 \\
    MSE (aug--clean) & 0.268 & 1 \\
    SIGReg & 2.334 & 1 \\
    \midrule
    Total & 2.838 & \\
    \bottomrule
  \end{tabular}
\end{table}

\subsection{Discussion}
\label{sec:discussion}

Within the fixed 3M architecture, \emph{the choice of SSL objective matters
more than incremental architectural changes}.
Dual InfoNCE and uniformity occupy opposite corners of the k-NN/linear
frontier (Figure~\ref{fig:tradeoff}); VICReg and SIGReg offer intermediate
and linear-favoring alternatives respectively.
Our asymmetric design explicitly composes instance contrast on the mask path
with spread on the augmentation path, yielding a practical knob for
low-resource CIFAR-100 training.

Outstanding work: complete the 200-epoch ablation table
(Table~\ref{tab:ablation-200}), finalize 500-epoch InfoNCE+SIGReg, and add
ResNet-50 scaling experiments (\texttt{latent\_infonce\_jepa\_sigreg\_resnet50}).
